\documentclass[11pt]{article}
\usepackage[a4paper,margin=1in]{geometry}
\usepackage{fontspec}
\usepackage[boldfont=false]{xeCJK}
\setCJKmainfont{FandolSong-Regular.otf}
\usepackage{amsmath}
\usepackage{amssymb}
\usepackage{array}
\usepackage{booktabs}
\usepackage{graphicx}
\usepackage{adjustbox}
\usepackage{enumitem}
\setlist[itemize]{topsep=2pt,partopsep=0pt,parsep=0pt,itemsep=2pt,leftmargin=1.4em}
\setlist[enumerate]{topsep=2pt,partopsep=0pt,parsep=0pt,itemsep=2pt,leftmargin=1.6em}
\usepackage{parskip}
\usepackage{fvextra}
\fvset{breaklines=true,breakanywhere=true,fontsize=\small}
\setlength{\parindent}{0pt}

\begin{document}

\begin{center}
  {\LARGE\bfseries Algorithm Design and Problem-solving}\\[0.35em]
  {\large A Level Computer Science}
\end{center}
\vspace{0.6em}

\section*{Computational thinking}

\textbf{Computational thinking} 计算思维 is the set of mental tools for \textbf{analysing a problem and designing a solution a computer can run}. Two key ones are abstraction and decomposition.

\subsection*{Abstraction}

\textbf{Abstraction} 抽象 means \textbf{keeping the essential features} of a problem and \textbf{ignoring the irrelevant detail}, giving a simpler model.

Examples: a train-network map keeps the stations and lines but drops geography; a class in object-oriented programming keeps the few attributes and methods the system needs; a function hides a piece of work behind a name. Abstraction is essential because a full model of any real problem would be too big to reason about.

\subsection*{Decomposition}

\textbf{Decomposition} 分解 means \textbf{breaking a large problem into smaller sub-problems}, each easier to solve and tackled one at a time.

\begin{enumerate}
\item find the main parts of the task.

\item break each into smaller sub-tasks.

\item continue until each is small enough to design directly.

\item solve the small tasks and combine them.

\end{enumerate}

For stock control: "manage stock" → "record sales", "record deliveries", "produce reports" → ("record sales") "look up product", "decrease stock count", "save the transaction". Decomposition makes big problems manageable, lets a team divide the work, and gives modular code.

\par\medskip
\noindent\begin{minipage}{\linewidth}
\centering
\adjustbox{max width=\linewidth, max totalheight=0.42\textheight}{%
  \includegraphics{/Users/joschka/Desktop/Programming/eXercise/output/eXam/handouts/a_level_computer_science/assets/09-p235-decomposition-modules.pdf}}\\[0.3em]
\emph{Decomposing a program into modules and sub-modules}
\end{minipage}\par\medskip

\section*{Algorithms}

An \textbf{algorithm} 算法 is a \textbf{solution expressed as a sequence of defined steps}. Each step is \textbf{unambiguous} 无歧义 (one meaning), \textbf{deterministic} 确定性 (same input → same output), \textbf{finite} (the steps end), and \textbf{effective} (each can be done). An algorithm says \emph{what to do}, independent of the programming language used to implement it.

\section*{Identifier table}

When you start an algorithm, list every piece of data in an \textbf{identifier table} 标识符表 — its \textbf{variable} 变量 name, \textbf{data type} 数据类型, and description:

\begin{center}
\adjustbox{max width=\linewidth}{%
\begin{tabular}{lll}
\toprule
\textbf{Variable name} & \textbf{Data type} & \textbf{Description} \\
\midrule
\texttt{Category} & \texttt{STRING} & the product category \\
\texttt{SaleDate} & \texttt{DATE} & when the item was sold \\
\texttt{ItemCost} & \texttt{REAL} & cost of the item \\
\texttt{InStock} & \texttt{BOOLEAN} & \texttt{TRUE} if in stock \\
\bottomrule
\end{tabular}}
\end{center}

Use \textbf{descriptive} names (\texttt{ItemCost}, not \texttt{x}); common types are \texttt{INTEGER}, \texttt{REAL}, \texttt{STRING}, \texttt{CHAR}, \texttt{BOOLEAN}, \texttt{DATE}, plus arrays. The table forces you to name every piece of data before writing code.

\section*{Pseudocode — the three basic constructs}

\textbf{Pseudocode} 伪代码 is a structured, language-neutral way to describe algorithms.

\subsection*{1. Sequence}

Steps run one after another (\textbf{sequence} 顺序):

\begin{Verbatim}
INPUT Name
INPUT Age
OUTPUT "Hello", Name
\end{Verbatim}

\subsection*{2. Selection}

A choice of which steps run, based on a condition (\textbf{selection} 选择):

\begin{Verbatim}
IF Age >= 18 THEN
    OUTPUT "Adult"
ELSE
    OUTPUT "Minor"
ENDIF
\end{Verbatim}

For more options, use \texttt{CASE OF ... ENDCASE}.

\subsection*{3. Iteration}

Repeating a block (\textbf{iteration} 迭代, a \textbf{loop} 循环):

\begin{Verbatim}
FOR i ← 1 TO 10
    OUTPUT i
NEXT i
\end{Verbatim}

A \textbf{WHILE} loop tests the condition \textbf{before} each pass (may run zero times); a \textbf{REPEAT...UNTIL} loop tests \textbf{after} each pass (always runs at least once).

\subsection*{Common operations}

\begin{itemize}
\item \textbf{assignment} 赋值: \texttt{x ← 5} (an arrow; \texttt{=} is for comparison).

\item input/output: \texttt{INPUT variable}, \texttt{OUTPUT expression}.

\item comparisons \texttt{=}, \texttt{\textless{}\textgreater{}}, \texttt{\textless{}}, \texttt{\textgreater{}}, \texttt{\textless{}=}, \texttt{\textgreater{}=}; logic \texttt{AND}, \texttt{OR}, \texttt{NOT}.

\item arithmetic \texttt{+ - * /}, plus \texttt{DIV} (integer division) and \texttt{MOD} (remainder).

\item strings: \texttt{LENGTH}, \texttt{LEFT}, \texttt{RIGHT}, \texttt{MID}, and \texttt{\&} for \textbf{concatenation} 拼接 (joining).

\end{itemize}

\subsection*{Input → Process → Output}

Every program follows this shape:

\begin{Verbatim}
INPUT Length
INPUT Width
Area ← Length * Width
OUTPUT "Area = ", Area
\end{Verbatim}

Listing the inputs and outputs first makes the algorithm cleaner.

\section*{Three notations}

The same algorithm can be written three ways.

\begin{itemize}
\item \textbf{structured English} 结构化英语 — natural language with indentation and fixed keywords; good for a high-level description.

\item \textbf{flowchart} 流程图 — a diagram with standard shapes:

\end{itemize}

\begin{center}
\adjustbox{max width=\linewidth}{%
\begin{tabular}{ll}
\toprule
\textbf{Shape} & \textbf{Meaning} \\
\midrule
Rounded rectangle & Start / Stop \\
Parallelogram & Input / Output \\
Rectangle & Process \\
Diamond & Decision \\
Arrow & Flow of control \\
\bottomrule
\end{tabular}}
\end{center}

\begin{itemize}
\item \textbf{pseudocode} — the keyword notation above; closest to code.

\end{itemize}

You should be able to convert between any pair: each \texttt{IF} is a decision diamond, each loop is a back-arrow, and a sequence is stacked rectangles.

\par\medskip
\noindent\begin{minipage}{\linewidth}
\centering
\adjustbox{max width=\linewidth, max totalheight=0.42\textheight}{%
  \includegraphics{/Users/joschka/Desktop/Programming/eXercise/output/eXam/handouts/a_level_computer_science/assets/09-p237-flowchart-example.pdf}}\\[0.3em]
\emph{A flowchart for averaging a list of numbers, using the standard shapes}
\end{minipage}\par\medskip

\section*{Stepwise refinement}

\textbf{Stepwise refinement} 逐步求精 starts with a high-level outline and \textbf{expands each step} until it is small enough to code. For an average of $n$ numbers:

\textbf{Level 1}:

\begin{Verbatim}
Read in the numbers
Compute the average
Output the average
\end{Verbatim}

\textbf{Level 2}:

\begin{Verbatim}
INPUT n
total ← 0
FOR i ← 1 TO n
    INPUT value
    total ← total + value
NEXT i
average ← total / n
OUTPUT average
\end{Verbatim}

Each refinement keeps the previous structure and adds detail.

\section*{Logic statements}

A \textbf{logic statement} 逻辑语句 is a \textbf{Boolean} 布尔 condition that controls branching, built from comparisons (\texttt{x \textgreater{} 10}), connectives (\texttt{AND}, \texttt{OR}, \texttt{NOT}) and brackets. Use it as the condition of \texttt{IF}, \texttt{WHILE} or \texttt{REPEAT...UNTIL}:

\begin{Verbatim}
WHILE attempts < 3 AND NOT loggedIn DO
    INPUT password
    IF password = correctPassword THEN
        loggedIn ← TRUE
    ELSE
        attempts ← attempts + 1
    ENDIF
ENDWHILE
\end{Verbatim}

\textbf{Precedence} 优先级 (highest to lowest): \texttt{NOT}, then \texttt{AND}, then \texttt{OR}. Use brackets when unsure. Common mistakes:

\begin{itemize}
\item \texttt{a = 1 OR 2} is wrong — write \texttt{a = 1 OR a = 2}.

\item \texttt{NOT a \textgreater{} 5} means \texttt{NOT (a \textgreater{} 5)}, i.e. \texttt{a \textless{}= 5}.

\item \texttt{NOT (A AND B)} is the same as \texttt{(NOT A) OR (NOT B)} (\textbf{De Morgan's law} 德摩根定律) — handy for simplifying conditions.

\end{itemize}


\end{document}
